\documentclass[../../main.tex]{subfiles}

\begin{document}

\section{Points and Circles of a Triangle}

There are so many important points circles and circles in triangles.
Knowing relationship between them can significantly help in solving
many problems. Before we get into the first content, let's clear
up some symbols. To be honest, I just don't think it would be
efficient to define each points, so I will just mention them here.
We will be using $O$ for circumcenter, $I$ for incenter, $H$ for
orthocenter, $G$ for centroid, and $I_A, I_B, I_C$ for excenters
unless otherwise stated. With that in mind, let's get started with
Euler line and related theorem!

\subsection{Servoir's Theorem and Euler Line}

As a warm-up, let's start with Servoir's Theorem and Euler Line.
Fun fact, when you search Servoir's Theorem online, you won't
probably get the theorem that we will discuss. I am not sure if
we have an English name for this theorem or if it is even named,
but this is actually very useful and Koreans call this Servoir's
Theorem.

\begin{theorem}{Servoir's Theorem}{Servoir's Theorem}
For triangle $ABC$ and midpoint $M$ on side $BC$, $AH = 2OM$.
\end{theorem}
\begin{proof}
First, consider the following diagram.

\begin{figure}[H]
\centering
\begin{tikzpicture}[scale=1.3]
    \coordinate (A) at (1, 4);
    \coordinate (B) at (0, 0);
    \coordinate (C) at (5, 0);
    \coordinate (D) at (1, 0);
    \coordinate (E) at (0.29411764705882354,1.1764705882352942);
    \coordinate (F) at (5, 3);
    \coordinate (M) at (2.5, 0);
    \coordinate (O) at (2.5, 1.5);
    \coordinate (H) at (1, 1);

    \draw[thick] (A) -- (B) -- (C) -- cycle;
    \draw[thick] (O) circle (2.91547594742265);
    \draw[thick, dotted] (A) -- (D);
    \draw[thick, dotted] (C) -- (E);
    \draw[thick, dotted, BrickRed] (A) -- (H);
    \draw[thick, dotted, BrickRed] (O) -- (M);
    \draw[thick, dotted] (A) -- (F);
    \draw[thick, dotted] (B) -- (F);
    \draw[thick, dotted] (C) -- (F);
    \draw[thick, dotted] (H) -- (C);
    \draw[fill=BrickRed,opacity=0.2]
        (A) -- (H) -- (C) -- (F) -- cycle;

    \pic [draw, angle radius=3mm] {right angle=O--M--C};
    \pic [draw, angle radius=3mm] {right angle=A--D--C};
    \pic [draw, angle radius=3mm] {right angle=B--A--F};
    \pic [draw, angle radius=3mm] {right angle=B--C--F};
    \pic [draw, angle radius=3mm] {right angle=C--E--B};

    \draw[fill=black] (A) circle (0.05)
        node[above left] {\small $A$};
    \draw[fill=black] (B) circle (0.05)
        node[below left] {\small $B$};
    \draw[fill=black] (C) circle (0.05)
        node[below right] {\small $C$};
    \draw[fill=black] (D) circle (0.05);
    \draw[fill=black] (E) circle (0.05);
    \draw[fill=black] (F) circle (0.05)
        node[above right] {\small $D$};
    \draw[fill=BrickRed] (M) circle (0.05)
        node[below] {\small $M$};
    \draw[fill=BrickRed] (O) circle (0.05)
        node[below right] {\small $O$};
    \draw[fill=BrickRed] (H) circle (0.05)
        node[above right] {\small $H$};
\end{tikzpicture}
\end{figure}

From the diagram above, it is evident that lines $\overline{AH}$
and $\overline{DC}$ are parallel. Similarly, $\overline{CH}$ and
$\overline{AD}$ are parallel. Therefore, quadrilateral $AHCD$ is
a parallelogram and $AH = CD = 2OM$.
\end{proof}

Continuing from the diagram above, we could also consider $G$ of
$\triangle{ABC}$.

\begin{figure}[H]
\centering
\begin{tikzpicture}[scale=1.3]
    \coordinate (A) at (1, 4);
    \coordinate (B) at (0, 0);
    \coordinate (C) at (5, 0);
    \coordinate (D) at (1, 0);
    \coordinate (E) at (0.29411764705882354,1.1764705882352942);
    \coordinate (F) at (5, 3);
    \coordinate (M) at (2.5, 0);
    \coordinate (O) at (2.5, 1.5);
    \coordinate (H) at (1, 1);
    \coordinate (G) at (2,1.3333333333333333);

    \draw[thick] (A) -- (B) -- (C) -- cycle;
    \draw[thick] (O) circle (2.91547594742265);
    \draw[thick, dotted] (A) -- (D);
    \draw[thick, dotted] (C) -- (E);
    \draw[thick, dotted, BrickRed] (A) -- (H);
    \draw[thick, dotted, BrickRed] (O) -- (M);
    \draw[thick, dotted] (A) -- (F);
    \draw[thick, dotted] (B) -- (F);
    \draw[thick, dotted] (C) -- (F);
    \draw[thick, dotted] (H) -- (C);
    \draw[thick] (A) -- (M);
    \draw[thick, BrickRed] (H) -- (G);
    \draw[thick, RoyalPurple] (G) -- (O);
    \draw[fill=BrickRed,opacity=0.2] (A) -- (H) -- (G) -- cycle;
    \draw[fill=RoyalPurple,opacity=0.2] (G) -- (O) -- (M) -- cycle;

    \pic [draw, angle radius=3mm] {right angle=O--M--C};
    \pic [draw, angle radius=3mm] {right angle=A--D--C};
    \pic [draw, angle radius=3mm] {right angle=B--A--F};
    \pic [draw, angle radius=3mm] {right angle=B--C--F};
    \pic [draw, angle radius=3mm] {right angle=C--E--B};

    \draw[fill=black] (A) circle (0.05)
        node[above left] {\small $A$};
    \draw[fill=black] (B) circle (0.05)
        node[below left] {\small $B$};
    \draw[fill=black] (C) circle (0.05)
        node[below right] {\small $C$};
    \draw[fill=black] (D) circle (0.05);
    \draw[fill=black] (E) circle (0.05);
    \draw[fill=black] (F) circle (0.05)
        node[above right] {\small $D$};
    \draw[fill=BrickRed] (M) circle (0.05)
        node[below] {\small $M$};
    \draw[fill=BrickRed] (O) circle (0.05)
        node[below right] {\small $O$};
    \draw[fill=BrickRed] (H) circle (0.05)
        node[above right] {\small $H$};
    \draw[fill=BrickRed] (G) circle (0.05)
        node[above left] {\small $G$};
\end{tikzpicture}
\end{figure}

Consider triangles $AHG$ and $MOG$. Notice that $\angle(HAG) =
\angle(OMG)$ since $\overline{AH} \parallel \overline{OM}$.
Moreover by Servoir's Theorem, $AH : OM = 2 : 1$ holds and
$AG : GM = 2 : 1$ is true by definition. Therefore, $\triangle{AHG}$
and $\triangle{MOG}$ are similar and $\angle(AGH) = \angle(MGO)$.
In other words, $H, G, O$ are collinear, and we have a special name
for the line!

\begin{definition}{}{}
The \textbf{Euler Line} of a triangle $ABC$ is the line that contains
$H, G, O$.
\end{definition}

As you can see, Euler line exists for all triangle except equilateral
triangle. The fact that $H, G, O$ is collinear will come very handy
when solving problems!

\subsection{Nine-Point Circle}



\end{document}


